\section{Introducción}
Los modelos contextuales representan una palabra mediante un vector que depende de la secuencia en la que aparece. Esta propiedad permite estudiar la variación de uso que una representación estática comprime en un único vector por entrada léxica. BERT constituye una referencia para esta familia de modelos \parencite{devlin2019}. Sin embargo, que dos ocurrencias tengan representaciones diferentes no establece que la diferencia corresponda a sentidos léxicos distintos. La contextualización también puede codificar variación sintáctica, temática o discursiva.

El estudio de la polisemia requiere, por tanto, relacionar una estructura geométrica continua con un inventario discreto de sentidos. Aquí se emplean las anotaciones WordNet de SemCor como referencia operacional: una vecindad será semánticamente homogénea cuando sus ocurrencias compartan synset. Esta decisión permite medir correspondencia con un recurso léxico, sin equiparar ese inventario con una partición exhaustiva del significado.

La dimensión elevada plantea un segundo problema. Una representación de \Num{d} coordenadas puede proyectarse al plano para facilitar su inspección, pero esa transformación modifica distancias y relaciones de vecindad. PCA privilegia la varianza, ISOMAP utiliza caminos en un grafo de proximidad y UMAP optimiza relaciones locales ponderadas. Sus objetivos no coinciden con maximizar la coincidencia de sentidos. En consecuencia, una figura con regiones de colores separadas puede resultar semánticamente sugerente y, al mismo tiempo, reemplazar muchos vecinos originales.

El análisis geométrico de representaciones contextualizadas tiene antecedentes que muestran la pertinencia de examinar sus propiedades espaciales \parencite{ethayarajh2019}. La evaluación en desambiguación léxica ofrece otra perspectiva sobre la información de sentido \parencite{hadiwinoto2019}. Nuestro interés específico es articular tres preguntas empíricas que suelen confundirse al interpretar proyecciones: cuántos vecinos sobreviven, qué tan alejados estaban los vecinos introducidos y cuántos vecinos comparten la anotación semántica.

La pregunta central es \emph{¿las relaciones de vecindad que conservan o introducen PCA, ISOMAP y UMAP corresponden a ocurrencias que comparten el mismo synset de WordNet?} Para responderla se comparan las identidades de los vecinos en diferentes escalas, se evalúa el cambio de pureza y selectividad semántica al cambiar la configuración, y se deriva una descomposición exacta del cambio de pureza en volumen de sustituciones y calidad semántica de los vecinos introducidos y eliminados. El análisis se circunscribe a los grupos observados y no pretende establecer un orden universal entre reductores.
